% =====================================================================================
% phan1-nhom.tex - PHẦN I: nội dung nhóm (giống nhau ở cả 4 báo cáo)
% =====================================================================================
\PhanLon{I. NỘI DUNG NHÓM}

\section{Mô tả kiến trúc chung của hệ thống}

\subsection{Bài toán và phạm vi}
\textbf{Ironfront Reborn} chuyển mã nguồn trò chơi bắn súng góc nhìn thứ nhất \textit{Ravenfield Beta 5}
(Unity 6000.3, vốn chỉ chơi đơn) thành một trò chơi \textbf{nhiều người chơi, máy chủ nắm quyền quyết định}
(server-authoritative). Yêu cầu cốt lõi của môn Lập trình mạng: \textbf{toàn bộ tầng mạng được viết tay trên
socket TCP và UDP thô} của hệ điều hành - không dùng Mirror, Photon, Netcode for GameObjects,
KCP, gRPC hay WebSocket.

\begin{center}\small
\begin{tabularx}{\linewidth}{|L{5.2cm}|Y|}
\hline
\hd{Tiêu chí} & \hd{Giá trị thiết kế} \\ \hline
Quy mô một trận & 16 người chơi thật + 32 bot AI (tối đa 64 actor, 16 phương tiện) \\ \hline
Tần số mô phỏng / snapshot / gửi input & 30 Hz / 20 Hz / 30 Hz \\ \hline
Băng thông xuống mỗi client & $\le$ 8 KB/s (đo trên server đã triển khai: 1,67 KB/s) \\ \hline
Triển khai & LAN trước, sau đó VPS công cộng (Docker Compose, TLS cho kênh TCP) \\ \hline
Công nghệ & .NET 8 (master server, công cụ), netstandard2.1 (thư viện dùng chung), Unity 6 (client và game server headless) \\ \hline
\end{tabularx}
\end{center}

\subsection{Kiến trúc ba tiến trình}
Hệ thống gồm \textbf{ba loại tiến trình} độc lập, giao tiếp qua \textbf{hai giao thức tự định nghĩa}:
\textbf{MSP} (Master Server Protocol, trên TCP) cho dịch vụ sảnh chờ và \textbf{GSP} (Game Server Protocol,
trên UDP) cho lưu lượng trong trận. Hai giao thức gặp nhau ở đúng một điểm: \textbf{joinTicket} được master
ký HMAC và game server tự kiểm tra.

\begin{figure}[H]
\centering
\begin{tikzpicture}[x=1cm, y=1cm]
  % ---- Client
  \node[bxB, text width=3.2cm] (cg) at (2,0)    {Gameplay\\{\footnotesize Actor, Weapon, UI}};
  \node[bxB, text width=3.2cm] (cp) at (2,-1.4) {Dự đoán + Nội suy};
  \node[bxB, text width=3.2cm] (cr) at (2,-2.6) {Replication Client};
  \node[bxB, text width=3.2cm] (ct) at (2,-3.8) {Transport UDP};
  \node[bxB, text width=3.2cm] (cm) at (2,-5.0) {MasterClient TCP};
  \draw[arr] (cg) -- (cp); \draw[arr] (cp) -- (cr); \draw[arr] (cr) -- (ct);
  \node[grp, fit=(cg)(cm)] (gC) {};
  \node[grpt] at (gC.north west) {Client - Unity 6};
  % ---- Game server
  \node[bxG, text width=3.4cm] (ss) at (14,-1.0) {Mô phỏng có thẩm quyền\\{\footnotesize PhysX, A*, AI gốc}};
  \node[bxG, text width=3.4cm] (sr) at (14,-2.4) {Replication Server\\{\footnotesize snapshot, delta, interest}};
  \node[bxG, text width=3.4cm] (st) at (14,-3.8) {Transport UDP};
  \node[bxG, text width=3.4cm] (sm) at (14,-5.0) {MasterLink TCP};
  \draw[arr] (st) -- (sr); \draw[arr] (sr) -- (ss);
  \node[grp, fit=(ss)(sm)] (gS) {};
  \node[grpt] at (gS.north west) {Game Server - Unity headless};
  % ---- UDP link
  \draw[biarr] (ct.east) -- node[lbl, above=2pt]{UDP $\cdot$ GSP $\cdot$ 20-30 Hz} (st.west);
  % ---- Master server
  \node[bxO, text width=4.2cm] (ml) at (8,-7.2) {Sảnh chờ, danh sách phòng};
  \node[bxO, text width=2.1cm] (ma) at (5.6,-8.8) {Xác thực,\\tài khoản};
  \node[bxO, text width=2.1cm] (mm) at (8,-8.8) {Ghép trận};
  \node[bxO, text width=2.1cm] (mg) at (10.4,-8.8) {Registry\\game server};
  \node[db, text width=1.6cm] (md) at (5.6,-10.3) {SQLite};
  \draw[arr] (ml) -- (ma); \draw[arr] (ml) -- (mm); \draw[arr] (ml) -- (mg); \draw[arr] (ma) -- (md);
  \node[grp, fit=(ml)(ma)(mg)(md)] (gM) {};
  \node[grpt] at (gM.north west) {Master Server - .NET 8};
  % ---- TCP links
  \draw[biarr] (cm.south) |- node[lbl, pos=0.75]{TCP $\cdot$ MSP\\đăng nhập, phòng, chat} (ml.west);
  \draw[biarr] (sm.south) |- node[lbl, pos=0.75]{TCP $\cdot$ MSP\\đăng ký, heartbeat} (ml.east);
\end{tikzpicture}
\caption{Kiến trúc tổng thể ba tiến trình và hai giao thức MSP/GSP}
\label{fig:n-kientruc}
\end{figure}

\begin{center}\small
\begin{tabularx}{\linewidth}{|L{2.3cm}|L{3.2cm}|Y|L{3.4cm}|}
\hline
\hd{Tiến trình} & \hd{Nền tảng} & \hd{Trách nhiệm chính} & \hd{Giao tiếp} \\ \hline
\textbf{Client} & Unity 6, bản build thường & Hiển thị, nhận input, dự đoán chuyển động người chơi cục bộ, nội suy actor ở xa, giao diện sảnh và HUD & TCP tới Master; UDP tới Game Server \\ \hline
\textbf{Game Server} & Unity 6 headless (chế độ không đồ hoạ) & Chạy mô phỏng có thẩm quyền (vật lý, AI, bắn, điểm chiếm), tạo snapshot riêng cho từng client, bù trễ & Nhận UDP từ client; TCP tới Master \\ \hline
\textbf{Master Server} & .NET 8 console, SQLite & Tài khoản, phiên, phòng, ghép trận, chat sảnh, registry game server, cấp joinTicket & Nhận TCP từ client và game server \\ \hline
\end{tabularx}
\end{center}

\subsection{Phân chia TCP và UDP theo đặc tính dữ liệu}
\begin{center}\small
\begin{tabularx}{\linewidth}{|L{2.8cm}|Y|Y|}
\hline
\hd{Đặc tính} & \hd{TCP - Master Server} & \hd{UDP - Game Server} \\ \hline
Dùng cho & Đăng nhập, danh sách phòng, tạo/vào phòng, ghép trận, chat sảnh, heartbeat server & Mọi lưu lượng trong trận: input, snapshot, sự kiện gameplay \\ \hline
Tần suất & Vài gói/phút mỗi client & 20-60 gói/giây mỗi chiều \\ \hline
Mất gói & Không chấp nhận (mất \code{LOGIN_RES} là treo màn hình) & Chấp nhận với trạng thái: snapshot sau thay thế snapshot trước \\ \hline
Nếu đổi chỗ & Phải tự viết lại độ tin cậy mà TCP đã làm tốt & Head-of-line blocking: một gói mất chặn mọi gói sau; Nagle + delayed ACK thêm 40-200 ms \\ \hline
Đóng khung & Tiền tố độ dài 4 byte + JSON & Datagram tự nhiên, header nhị phân 16 byte \\ \hline
\end{tabularx}
\end{center}
Không dùng WebSocket: WebSocket chạy trên TCP nên thừa hưởng nguyên vẹn head-of-line blocking, cộng thêm chi
phí bắt tay HTTP - một ràng buộc chỉ cần cho trình duyệt, không cần cho client desktop.

\subsection{Kiến trúc phân lớp hệ thống}
Mã nguồn chia thành các thư viện dùng chung độc lập hoàn toàn với engine đồ hoạ; dự án Unity tích hợp chúng dưới dạng các thư viện liên kết động (DLL). Nhờ vậy logic mạng được kiểm thử tự động độc lập nhanh chóng mà không cần mở Unity Editor.

\begin{figure}[H]
\centering
\begin{tikzpicture}[x=1cm, y=1cm, lay/.style={box, text width=3.35cm, minimum height=1.25cm},
                    rowname/.style={font=\small\bfseries, align=left, text width=2.4cm, anchor=west}]
  \node[rowname] at (0,0) {Ứng dụng};
  \node[lay, fill=cOrange] (u1) at (4.6,0)  {\textbf{Trò chơi Unity}\\{\footnotesize client + server headless}};
  \node[lay, fill=cPurple] (u2) at (8.35,0) {\textbf{Master Server}\\{\footnotesize TCP, CSDL SQLite}};
  \node[lay, fill=cPurple] (u3) at (12.1,0) {\textbf{Công cụ phụ trợ}\\{\footnotesize Kiểm thử tải, phát lại gói}};
  \node[rowname] at (0,-2.1) {Thư viện mạng\\{\footnotesize\mdseries dùng chung}};
  \node[lay, fill=cBlue, text width=2.55cm] (l1) at (4.2,-2.1)  {\textbf{Tầng Giao vận}\\{\footnotesize UDP tin cậy}};
  \node[lay, fill=cBlue, text width=2.55cm] (l2) at (7.1,-2.1)  {\textbf{Tầng Đồng bộ}\\{\footnotesize snapshot, bù trễ}};
  \node[lay, fill=cBlue, text width=2.55cm] (l3) at (10.0,-2.1) {\textbf{Kết nối Master}\\{\footnotesize client MSP}};
  \node[lay, fill=cBlue, text width=2.55cm] (l4) at (12.9,-2.1) {\textbf{Liên kết GS}\\{\footnotesize GS báo master}};
  \node[rowname] at (0,-4.2) {Tầng nền};
  \node[lay, fill=cGreen, text width=6.2cm] (p1) at (6.05,-4.2) {\textbf{Giao thức chung (Protocol)}\\{\footnotesize Hằng số, header, cấu trúc gói, lượng tử hoá, joinTicket}};
  \node[lay, fill=cGreen, text width=4.3cm] (p2) at (11.95,-4.2) {\textbf{Cấu hình chung}\\{\footnotesize Biến môi trường hệ thống}};
  \node[rowname] at (0,-6.3) {Bảo đảm\\chất lượng};
  \node[lay, fill=cGray, text width=6.2cm] (q1) at (6.05,-6.3) {\textbf{Kiểm thử tự động}\\{\footnotesize Kiểm thử giao thức, giao vận, đồng bộ, máy chủ}};
  \node[lay, fill=cGray, text width=4.3cm] (q2) at (11.95,-6.3) {\textbf{Bộ kiểm tra đặc tả}\\{\footnotesize chặn build khi lệch đặc tả}};
  \draw[arr] (4.2,-1.4) -- (4.2,-0.65);
  \draw[arr] (10.0,-1.4) -- (10.0,-0.65);
  \node[lbl, anchor=west] at (4.35,-1.02) {Tích hợp DLL};
  \node[lbl, anchor=west] at (10.15,-1.02) {tham chiếu trực tiếp};
  \draw[arr] (4.2,-3.55) -- (4.2,-2.75);
  \draw[arr] (7.1,-3.55) -- (7.1,-2.75);
  \draw[arr] (11.95,-3.55) -- (11.95,-2.75);
  \draw[arr] (3.6,-5.65) -- (3.6,-4.85);
  \draw[cLine, dashed] (2.9,-5.25) -- (14.2,-5.25);
\end{tikzpicture}
\caption{Phân lớp kiến trúc và quan hệ phụ thuộc giữa các phân hệ}
\label{fig:n-phanlop}
\end{figure}

\begin{center}\small
\begin{tabularx}{\linewidth}{|L{4.6cm}|L{2.4cm}|L{1.5cm}|Y|}
\hline
\hd{Phân hệ / Thành phần} & \hd{Loại} & \hd{Quy mô} & \hd{Vai trò nghiệp vụ} \\ \hline
Giao thức chung (Protocol) & Thư viện dùng chung & 5\,300 dòng & Định dạng gói tin, chuẩn tuần tự hoá và hằng số giao thức \\ \hline
Tầng Giao vận (Transport) & Thư viện giao vận & 4\,100 dòng & Độ tin cậy trên UDP, kênh truyền, phân mảnh, kiểm soát nghẽn \\ \hline
Đồng bộ và Mô phỏng & Thư viện logic & 18\,500 dòng & Snapshot, mã hóa delta, vùng quan tâm, tick server, bù trễ \\ \hline
Máy chủ trung tâm (Master) & Ứng dụng dịch vụ & 4\,300 dòng & Tiếp nhận TCP, xác thực, quản lý phòng, ghép trận, bảo mật \\ \hline
Liên kết máy chủ trung tâm & Thư viện kết nối & 1\,250 dòng & API trao đổi thông tin với máy chủ trung tâm \\ \hline
Tích hợp giao diện client & Tích hợp ứng dụng & 14\,600 dòng & Điều phối luồng, dự đoán, nội suy, giao diện sảnh và HUD \\ \hline
\end{tabularx}
\end{center}

\subsection{Kiến trúc triển khai}
\begin{figure}[H]
\centering
\begin{tikzpicture}[x=1cm, y=1cm]
  \node[bxB, text width=2.4cm] (c1) at (1.5,0)    {Client Unity 1};
  \node[bxB, text width=2.4cm] (c2) at (1.5,-2.4) {Client Unity N};
  \node[grp, fit=(c1)(c2)] (gP) {};
  \node[grpt] at (gP.north west) {Người chơi};
  \node[bxG, text width=2.8cm] (g1) at (6.3,0)    {gameserver-1\\{\footnotesize UDP 27015}};
  \node[bxG, text width=2.8cm] (g2) at (6.3,-2.4) {gameserver-2\\{\footnotesize UDP 27016}};
  \node[bxO, text width=3.2cm] (m)  at (6.3,-4.9) {master\\{\footnotesize TCP 27000 (TLS)\\metrics 27001 loopback}};
  \node[db, text width=1.9cm] (v)   at (13.4,-4.9) {volume\\SQLite};
  \node[bxY, text width=2.9cm] (a)  at (11.4,-2.4) {ironfront-alert\\{\footnotesize timer: đọc metrics}};
  \node[bxY, text width=2.9cm] (b)  at (11.4,0)    {ironfront-backup\\{\footnotesize timer: sao lưu}};
  \node[grp, fit=(g1)(m)(v)(a)(b)] (gV) {};
  \node[grpt] at (gV.north west) {Máy ảo VPS - Docker Compose};
  \node[bxP, text width=2.8cm] (reg) at (6.3,-7.6) {GHCR\\{\footnotesize image theo digest}};
  \draw[biarr] (c1) -- node[lbl, above=2pt]{UDP $\cdot$ GSP} (g1);
  \draw[biarr] (c2) -- node[lbl, above=2pt]{UDP $\cdot$ GSP} (g2);
  \draw[biarr] (gP.south) |- node[lbl, pos=0.72, above=2pt]{TCP/TLS $\cdot$ MSP} (m.west);
  \draw[arr] (g2.south) -- (m.north);
  \draw[arr] (g1.east) -- (8.6,0) -- node[lbl, sloped, above=1pt, pos=0.45]{đăng ký, heartbeat (TLS)} (8.6,-4.55) -- ($(m.east)+(0,0.35)$);
  \draw[cLine, thick] ($(m.east)+(0,-0.35)$) -- ($(v.west)+(0,-0.35)$);
  \draw[darr] (b.east) -- (13.4,0) -- (v.north);
  \draw[darr] (reg.north) -- node[lbl, right=2pt]{pull image} (reg.north |- gV.south);
\end{tikzpicture}
\caption{Sơ đồ triển khai trên một máy ảo: một master và hai game server headless}
\label{fig:n-trienkhai}
\end{figure}

Image được triển khai theo bản dựng cố định; khóa bí mật dùng chung và chứng chỉ TLS được cấu hình qua biến môi trường bảo mật độc lập. Sao lưu và cảnh báo chạy bằng bộ định thời hệ thống \textit{ngoài} container, vì bộ phát hiện
sự cố phải nằm ngoài thứ nó giám sát. Hạ tầng còn có cấu hình tự động hoá máy ảo và phương
án triển khai đám mây thay thế.

\subsection{Phân công theo phân hệ}
\begin{center}\small
\begin{tabularx}{\linewidth}{|L{3.8cm}|L{2.4cm}|L{4.2cm}|Y|}
\hline
\hd{Sinh viên} & \hd{Mã SV} & \hd{Phân hệ} & \hd{Phạm vi nhiệm vụ chính} \\ \hline
Nguyễn Thành Trung & B23DCCN861 & Client Unity & Giao diện sảnh, HUD, dự đoán và nội suy chuyển động client \\ \hline
Nguyễn Minh Toàn & B23DCCN832 & Tầng giao vận UDP & Giao thức UDP tin cậy, kênh truyền, phân mảnh, ACK bitfield \\ \hline
Nguyễn Duy Nghĩa & B23DCCN600 & Replication, mô phỏng server & Vòng tick mô phỏng, đồng bộ snapshot, nén delta và bù trễ \\ \hline
Nguyễn Tự Kiên & B23DCCN465 & Master Server TCP, hạ tầng & Máy chủ TCP, xác thực, quản lý phòng, ghép trận và hạ tầng VPS \\ \hline
\end{tabularx}
\end{center}

% -------------------------------------------------------------------------------------
\section{Mô tả thiết kế chung của hệ thống}

\subsection{Mô hình thẩm quyền}
Hệ thống \textbf{hoàn toàn do server quyết định}: client chỉ gửi \textit{ý định} (input), không bao giờ gửi kết
quả. Nhờ vậy gian lận tốc độ, dịch chuyển tức thời, bắn nhanh hay đạn vô hạn trở nên bất khả thi về cấu trúc
thay vì chỉ bị phát hiện.

\begin{center}\small
\begin{tabularx}{\linewidth}{|L{4.6cm}|L{2.3cm}|Y|}
\hline
\hd{Vấn đề} & \hd{Ai quyết định} & \hd{Client được làm gì} \\ \hline
Vị trí nhân vật & Server & Dự đoán ngay, sau đó hoà giải với server \\ \hline
Trúng đạn, sát thương & Server & Hiện hiệu ứng dự đoán, chờ \code{S_HIT_CONFIRM} \\ \hline
Máu, chết, hồi sinh & Server & Chỉ hiển thị \\ \hline
Độ tản đạn, giật súng & Server & Giật cục bộ cho cảm giác, không ảnh hưởng viên đạn thật \\ \hline
Bot AI, điểm chiếm, điểm số & Server & Nội suy, hiển thị \\ \hline
Lên xuống ghế phương tiện & Server & Gửi yêu cầu \code{C_SEAT_REQUEST} \\ \hline
Camera, UI, âm thanh, ragdoll & Client & Toàn quyền, không đồng bộ \\ \hline
\end{tabularx}
\end{center}

\subsection{Luồng nghiệp vụ đầu - cuối}
\begin{figure}[H]
\centering
\begin{tikzpicture}[x=1cm, y=1cm]
  \seqactor{C}{1.2}{Client}
  \seqactor{M}{7.8}{Master Server (TCP)}
  \seqactor{G}{14.4}{Game Server (UDP)}
  \seqbegin
  \seqmsg{G}{M}{GS\_REGISTER \{serverSecret, publicIp, udpPort\}}
  \seqret{M}{G}{GS\_REGISTER\_RES \{serverId\}}
  \seqmsg{C}{M}{LOGIN\_REQ \{username, SHA256(pass+user)\}}
  \seqret{M}{C}{LOGIN\_RES \{sessionToken, playerId\}}
  \seqmsg{C}{M}{ROOM\_LIST\_REQ}
  \seqret{M}{C}{ROOM\_LIST\_RES [phòng]}
  \seqmsg{C}{M}{ROOM\_JOIN\_REQ \{roomId\}}
  \seqself{M}{chọn game server còn chỗ; ký joinTicket =\\HMAC-SHA256(playerId, serverId, roomId, exp)}
  \seqret{M}{C}{ROOM\_JOIN\_RES \{ip, port, joinTicket\}}
  \seqmsg{C}{G}{CONNECT\_REQUEST \{version, joinTicket, clientSalt\}}
  \seqret{G}{C}{CONNECT\_CHALLENGE \{serverSalt\}}
  \seqmsg{C}{G}{CONNECT\_RESPONSE \{salt XOR, joinTicket\}}
  \seqself{G}{kiểm tra HMAC và hạn dùng của vé\\(không truy vấn lại Master Server)}
  \seqret{G}{C}{CONNECT\_ACCEPTED \{connectionId, serverTick\}}
  \seqret{G}{C}{S\_MATCH\_STATE + S\_SPAWN\_ACTOR $\times$ N}
  \seqloopbegin
  \seqmsg{C}{G}{C\_INPUT \{tick, các frame gần nhất\} - 30 Hz}
  \seqret{G}{C}{S\_SNAPSHOT \{tick, lastProcessedInputTick, delta\} - 20 Hz}
  \seqloopend{C}{G}{lặp}
  \seqmsg{C}{G}{C\_INPUT \{nút FIRE, yaw, pitch\}}
  \seqself{G}{tua hitbox về thời điểm client quan sát,\\raycast và phân xử trúng}
  \seqret{G}{C}{S\_HIT\_CONFIRM \{targetId, damage\}}
  \seqmsg{G}{M}{GS\_MATCH\_ENDED \{kết quả\}}
  \seqend{C,M,G}
\end{tikzpicture}
\caption{Luồng đầu - cuối: đăng nhập, vào phòng, kết nối game server, chơi, kết thúc trận}
\label{fig:n-luong}
\end{figure}

\subsection{Các kỹ thuật netcode cốt lõi}
\begin{figure}[H]
\centering
\begin{tikzpicture}[x=1cm, y=1cm]
  \seqactor{C}{2.5}{Client}
  \seqactor{S}{13}{Server}
  \seqbegin
  \seqself{C}{tick 100: nhấn W; dự đoán di chuyển ngay,\\lưu input tick 100}
  \seqmsg{C}{S}{C\_INPUT \{tick 100, moveZ +1\}}
  \seqself{S}{xử lý input tick 100,\\tính ra vị trí P}
  \seqret{S}{C}{S\_SNAPSHOT \{lastProcessedInputTick 100, pos P\}}
  \seqself{C}{Hoà giải: so vị trí dự đoán tại tick 100 với P;\\lệch quá ngưỡng thì đặt về P, phát lại input 101..hiện tại}
  \global\advance\seqy by -0.3cm
  \seqself{C}{actor ở xa: nội suy tại (now $-$ 100 ms)\\giữa hai snapshot đã nhận}
  \seqend{C,S}
\end{tikzpicture}
\caption{Dự đoán phía client, hoà giải với server và nội suy actor ở xa}
\label{fig:n-netcode}
\end{figure}

\begin{center}\small
\begin{tabularx}{\linewidth}{|L{4cm}|Y|L{4.2cm}|}
\hline
\hd{Kỹ thuật} & \hd{Mục đích} & \hd{Tham số} \\ \hline
Dự đoán phía client & Người chơi cục bộ di chuyển ngay khi nhấn phím & Cùng thuật toán mô phỏng chuyển động ở client và server \\ \hline
Hoà giải & Sửa sai khi dự đoán lệch, phát lại input chưa xác nhận & Ngưỡng lệch 0,25 m \\ \hline
Nội suy thực thể & Actor ở xa chuyển động mượt dù snapshot chỉ 20 Hz & Bộ đệm 100 ms \\ \hline
Bù trễ & Người ping cao vẫn bắn trúng thứ họ thấy & Tua tối đa 200 ms, lịch sử hitbox 1 s \\ \hline
\end{tabularx}
\end{center}

\subsection{Thiết kế giao thức}
\subsubsection{GSP - cấu trúc một datagram PAYLOAD}
\begin{figure}[H]
\centering
\begin{tikzpicture}[x=1cm, y=1cm, fld/.style={draw=cLine, minimum height=1.2cm, align=center, font=\footnotesize, inner sep=2pt}]
  \node[fld, fill=cBlue,   minimum width=4.6cm, anchor=west] (h) at (0,0)    {\textbf{GSP header}\\16 byte};
  \node[fld, fill=cGreen,  minimum width=2.2cm, anchor=west] (e) at (h.east) {\textbf{Envelope}\\3 byte};
  \node[fld, fill=cOrange, minimum width=2.8cm, anchor=west] (f) at (e.east) {\textbf{Frame}\\channelId,\\messageCount};
  \node[fld, fill=cPurple, minimum width=3.2cm, anchor=west] (m1) at (f.east) {\textbf{Message 1}\\msgType, length,\\body};
  \node[fld, fill=cPurple, minimum width=3.2cm, anchor=west] (m2) at (m1.east) {\textbf{Message n}\\\ldots};
  \draw[decorate, decoration={brace, amplitude=5pt, mirror}] (h.south west) -- node[below=6pt, font=\footnotesize]{byte 0--15} (h.south east);
  \draw[decorate, decoration={brace, amplitude=5pt, mirror}] (e.south west) -- node[below=6pt, font=\footnotesize]{16--18} (e.south east);
  \draw[decorate, decoration={brace, amplitude=5pt, mirror}] (f.south west) -- node[below=6pt, font=\footnotesize]{frame $\le$ 1181 byte; cả datagram $\le$ 1200 byte (MTU an toàn)} (m2.south east);
\end{tikzpicture}
\caption{Ba lớp của một datagram PAYLOAD: header GSP, envelope kênh của transport, frame message của ứng dụng}
\label{fig:n-datagram}
\end{figure}

\begin{center}\small
\begin{tabularx}{\linewidth}{|L{1.1cm}|L{4.2cm}|Y|L{3.6cm}|}
\hline
\hd{Kênh} & \hd{Kiểu} & \hd{Dữ liệu} & \hd{Khi mất gói} \\ \hline
0 & Không tin cậy, không thứ tự & Ping/pong & Bỏ qua \\ \hline
1 & Không tin cậy, có thứ tự (server $\rightarrow$ client) & \code{S_SNAPSHOT}, \code{S_WEAPON_FIRE}, \code{S_VEHICLE_SNAPSHOT} & Bỏ qua; gói cũ bị huỷ \\ \hline
2 & Tin cậy, có thứ tự & Sự kiện gameplay, spawn/despawn, chat, \code{C_ACK_BASELINE} & Gửi lại tới khi được ack \\ \hline
3 & Không tin cậy, có thứ tự (client $\rightarrow$ server) & \code{C_INPUT}, \code{C_VEHICLE_INPUT} & Bỏ qua; bù bằng dư thừa \\ \hline
\end{tabularx}
\end{center}
Độ tin cậy áp dụng \textbf{theo kênh}, không theo kết nối: một sự kiện bị mất trên kênh 2 không làm nghẽn snapshot
trên kênh 1 - ưu điểm cốt lõi so với TCP vốn ép mọi thứ vào một luồng.

\subsubsection{MSP - khung tin trên TCP}
\begin{figure}[H]
\centering
\begin{tikzpicture}[x=1cm, y=1cm, fld/.style={draw=cLine, minimum height=1.1cm, align=center, font=\footnotesize, inner sep=2pt}]
  \node[fld, fill=cBlue,   minimum width=3.6cm, anchor=west] (a) at (0,0)    {\textbf{length}\\u32, big-endian};
  \node[fld, fill=cGreen,  minimum width=3.2cm, anchor=west] (b) at (a.east) {\textbf{msgType}\\u16, little-endian};
  \node[fld, fill=cOrange, minimum width=6.0cm, anchor=west] (c) at (b.east) {\textbf{body}\\UTF-8 JSON, tối đa 64 KB};
\end{tikzpicture}
\caption{Khung tin MSP: tiền tố độ dài giải quyết bài toán ranh giới thông điệp trên luồng byte TCP}
\label{fig:n-msp}
\end{figure}

\subsubsection{joinTicket - cầu nối TCP và UDP}
\begin{figure}[H]
\centering
\begin{tikzpicture}[x=1cm, y=1cm, fld/.style={draw=cLine, minimum height=1.1cm, align=center, font=\footnotesize, inner sep=2pt}]
  \node[fld, fill=cBlue,   minimum width=2.0cm, anchor=west] (a) at (0,0)    {\textbf{playerId}\\u32};
  \node[fld, fill=cBlue,   minimum width=1.7cm, anchor=west] (b) at (a.east) {\textbf{serverId}\\u16};
  \node[fld, fill=cBlue,   minimum width=1.6cm, anchor=west] (c) at (b.east) {\textbf{roomId}\\u16};
  \node[fld, fill=cGreen,  minimum width=3.0cm, anchor=west] (d) at (c.east) {\textbf{expiresAtUnixMs}\\u64 (60 s)};
  \node[fld, fill=cOrange, minimum width=3.0cm, anchor=west] (e) at (d.east) {\textbf{displayName}\\16 byte UTF-8};
  \node[fld, fill=cPurple, minimum width=3.6cm, anchor=west] (f) at (e.east) {\textbf{hmac}\\HMAC-SHA256, 32 byte};
  \draw[decorate, decoration={brace, amplitude=5pt, mirror}] (a.south west) -- node[below=6pt, font=\footnotesize]{32 byte dữ liệu được ký} (e.south east);
  \draw[decorate, decoration={brace, amplitude=5pt, mirror}] (f.south west) -- node[below=6pt, font=\footnotesize]{32 byte chữ ký} (f.south east);
\end{tikzpicture}
\caption{Cấu trúc joinTicket 64 byte}
\label{fig:n-ticket}
\end{figure}
Master ký vé bằng khoá bí mật; game server kiểm tra chữ ký bằng thuật toán so sánh thời gian hằng số và kiểm hạn dùng.
Không cần gọi lại master, nên \textbf{master có thể dừng giữa trận mà trận vẫn tiếp tục}.

\subsection{Tham số thời gian}
\begin{center}\small
\begin{tabularx}{\linewidth}{|L{5.4cm}|L{3.2cm}|Y|}
\hline
\hd{Tham số} & \hd{Giá trị} & \hd{Ghi chú} \\ \hline
Tick mô phỏng server & 30 Hz (33,3 ms) & Theo đồng hồ thực, không theo chu kỳ cập nhật vật lý \\ \hline
Snapshot & 20 Hz (50 ms) & Rẻ hơn 30 Hz, nội suy bù độ mượt \\ \hline
Gửi input & 30 Hz & Mỗi gói lặp lại các frame gần nhất \\ \hline
Bộ đệm nội suy & 100 ms & Chịu được mất một snapshot \\ \hline
Cửa sổ bù trễ / lịch sử hitbox & 200 ms / 1 s & Chặn lạm dụng ping giả \\ \hline
Keep-alive / timeout UDP & 1 s / 10 s & \\ \hline
Heartbeat MSP client / game server & 15 s & Phát hiện kết nối TCP nửa mở \\ \hline
\end{tabularx}
\end{center}

\subsection{Phân loại đối tượng đồng bộ}
\begin{center}\small
\begin{tabularx}{\linewidth}{|L{3.4cm}|L{5.2cm}|Y|}
\hline
\hd{Loại} & \hd{Ví dụ} & \hd{Cách xử lý} \\ \hline
Actor được đồng bộ & Người chơi, bot & Có \code{actorId} u16, có trong snapshot, mã hoá delta \\ \hline
Phương tiện & Xe, trực thăng, tàu & Luồng riêng \code{S_VEHICLE_SNAPSHOT} \\ \hline
Chỉ ở server & Lịch sử hitbox, dữ liệu AI & Không bao giờ gửi \\ \hline
Chỉ ở client & Camera, UI, particle, ragdoll, âm thanh & Không bao giờ gửi \\ \hline
Sự kiện một lần & Chết, nổ, phát súng, chiếm điểm & Kênh tin cậy có thứ tự \\ \hline
Tĩnh & Địa hình, công trình, điểm hồi sinh & Có sẵn trong scene hai phía \\ \hline
\end{tabularx}
\end{center}

\subsection{Bảo mật ở mức phù hợp phạm vi}
\begin{center}\small
\begin{tabularx}{\linewidth}{|L{5cm}|Y|L{2.5cm}|}
\hline
\hd{Mối đe doạ} & \hd{Biện pháp} & \hd{Tầng} \\ \hline
Tăng tốc, dịch chuyển, bắn nhanh, đạn vô hạn & Server chỉ nhận input; kẹp tốc độ; tự quản lý cooldown và đạn & Replication \\ \hline
Mạo danh người chơi & \code{connectionId} gắn với \code{playerId} trong joinTicket đã ký & Transport \\ \hline
Khuếch đại giả mạo IP & Bắt tay challenge không lưu trạng thái (cookie HMAC) & Transport \\ \hline
Gói rác, lụt gói & Kiểm \code{protocolId}; giới hạn tốc độ theo IP; giới hạn nhóm phân mảnh & Transport \\ \hline
Nghe lén mật khẩu & TLS và ghim vân tay chứng chỉ; SHA-256 phía client, bcrypt phía server & Master \\ \hline
Tràn bộ nhớ, Slowloris & Khung tối đa 64 KB; hạn 30 s trước xác thực; 5 kết nối/IP & Master \\ \hline
SQL injection, dò mật khẩu & Truy vấn tham số hoá; 5 lần/phút/IP; khoá 15 phút sau 10 lần sai & Master \\ \hline
\end{tabularx}
\end{center}

\subsection{Các quyết định kiến trúc đã chốt}
\begin{center}\small
\begin{tabularx}{\linewidth}{|L{1.3cm}|Y|Y|}
\hline
\hd{Mã} & \hd{Quyết định} & \hd{Lý do} \\ \hline
AD-1 & Server nắm quyền, không có listen-server & Thẩm quyền rõ ràng; tránh sửa 21 singleton \\ \hline
AD-2 & Game server là Unity headless & Tái dùng PhysX, A*, AI; tiết kiệm 8-12 tuần \\ \hline
AD-3 & Không theo đuổi tất định (lockstep) & Bộ sinh ngẫu nhiên rải rác; engine vật lý không tất định \\ \hline
AD-4 & Ragdoll chỉ là hiệu ứng cục bộ & Tiết kiệm khoảng 1,7 MB/s \\ \hline
AD-5 & Transport, Replication là thư viện .NET thuần & Kiểm thử bằng xUnit, tránh xung đột Unity \\ \hline
AD-7 & Snapshot không tin cậy, sự kiện tin cậy & Đúng bản chất trạng thái và sự kiện \\ \hline
AD-8 & Không dùng WebSocket & Yêu cầu môn học và head-of-line blocking \\ \hline
\end{tabularx}
\end{center}

\subsection{Quy trình phát triển và bảo đảm chất lượng}
\begin{itemize}
  \item \textbf{Đặc tả giao thức đóng băng} là nguồn sự thật duy nhất của hệ thống; công cụ kiểm tra tự động sẽ cảnh báo nếu cấu trúc gói tin của các bên không đồng nhất.
  \item \textbf{Tách người thiết kế và người kiểm thử}: nhóm giao vận xây dựng cơ chế đóng gói dữ liệu, nhóm đồng bộ xây dựng bộ kiểm thử tính tương thích theo đặc tả.
  \item \textbf{Tích hợp liên tục (CI)} chạy tự động: kiểm tra toàn diện, coi cảnh báo là lỗi, chạy toàn bộ bộ kịch bản kiểm thử.
  \item \textbf{Tiêu chuẩn hoàn thành (Definition of Done)}: hoàn thành kiểm thử thực tế, toàn bộ kiểm nghiệm tự động đạt chuẩn, ứng dụng biên dịch sạch sẽ, tài liệu kỹ thuật hoàn thiện.
\end{itemize}
